%%%%%%%%%%%%%%%%%%%%%%%%%%%%%%%%%%%%%%%%%%%%%%%%
% E.Bigeard - emmanuel.bigeard@upsti.fr
% http://s2i.bigeard.me
% CC BY-NC-SA 2.0 FR - http://creativecommons.org/licenses/by-nc-sa/2.0/fr/
%%%%%%%%%%%%%%%%%%%%%%%%%%%%%%%%%%%%%%%%%%%%%%%%
\documentclass[11pt]{article}

%%%%%%%%%%%%%%%%%%%%%%%%%%%%%%%%%%%%%%%%%%%%%%%%
% Package UPSTI_Document
%%%%%%%%%%%%%%%%%%%%%%%%%%%%%%%%%%%%%%%%%%%%%%%% 
\usepackage{UPSTI_Document}

%---------------------------------%
% Paramètres du package
%---------------------------------%

% Version du document (pour la compilation)
% 1: Document prof
% 2: Document élève
% 3: Document à publier
\newcommand{\UPSTIidVersionDocument}{3}

% Affichage personnalisé de la classe
\newcommand{\UPSTIclasse}{SNT}

% Type de document
% 0: Custom*				7: Fiche Méthode			14: Document Réponses
% 1: Cours (par défaut)		8: Fiche Synthèse    		15: Programme de colle 
% 2: TD     				9: Formulaire
% 3: TP						10: Memo
% 4: Colle					11: Dossier Technique
% 5: DS						12: Dossier Ressource
% 6: DM						13: Concours Blanc
% * Si on met la valeur 0, il faut décommenter la ligne suivante: 		
%\newcommand{\UPSTItypeDocument}{Custom}
\newcommand{\UPSTIidTypeDocument}{3}

% Titre dans l'en-tête
\newcommand{\UPSTItitreEnTete}{LES PRODUITS ET LE PANIER}      
\newcommand{\UPSTItitreEnTetePages}{Site web marchand}      
\newcommand{\UPSTIsousTitreEnTete}{Site web marchand}      

% Version du document
\newcommand{\UPSTInumeroVersion}{1.0}
%v1.1 - Ajout de la chaîne fonctionnelle

%----------------------------------------------- 
\UPSTIcompileVars		% "Compile" les variables
%%%%%%%%%%%%%%%%%%%%%%%%%%%%%%%%%%%%%%%%%%%%%%%% 


%%%%%%%%%%%%%%%%%%%%%%%%%%%%%%%%%%%%%%%%%%%%%%%% 
% Début du document
%%%%%%%%%%%%%%%%%%%%%%%%%%%%%%%%%%%%%%%%%%%%%%%% 
\begin{document}
\UPSTIbuildPage

\UPSTIobjectif{Sur cette activité vous allez devoir réaliser différentes actions afin de faire fonctionner le coeur de votre site web marchand, c'est à dire la liste de vos produits 
et le panier d'achat. Le travail à réaliser se répartie en 3 rôles détaillé ci-dessous.
}

\competencetable{Compétences à valider}{
\competence{Savoir organiser les fichiers et les dossiers dans l'espace réseau personnel}
\competence{Définir une donnée personnelle.}
}

%Actions pour visualiser les produits et le panier d'achat
% - Mettre en place le serveur web
% - Télécharger les fichiers depuis Moodle
% - Ajouter les scripts nécessaires au fonctionnement du panier et à l'affichage des produits
% - Ajouter les images des produits
% - Compléter le fichier produit yaml

\section*{Rôle 1 - 2 Elèves : Mise en place du serveur web}

\activiteressources
{Mettre en place sur la page produit.}
{
\task Démarrer le raspberry PI et copier les fichiers du site web disponible sur Moodle ;
\task Ajouter au page "produits.html" et produit_detail.html" les scripts nécessaires au fonctionnement du panier et à l'affichage des produits ;
\task Lancer le serveur web à l'aide du programme Python ;
\task Tester le site Web ;
\task Attribuer dans le serveur DNS l'adresse web pour votre site marchand ;
\task Créer un point d'accès wifi ;
\task Tester l'accès au site web marchand avec votre smartphone.
}
{[RES]_Mettre en place un serveur web}
{Moodle}

\section*{Rôle 2 : 2 Elèves : Gestion des produits}

\activiteressources
{Compléter le fichier produits.yaml afin d'y ajouter vos propre produit.}
{
\task Ajouter les produits en respectant la structure du fichier YAML ;
\task Rechercher des images libres de droits pour illustrer vos produits ;
\task Ajouter les images dans le dossier "images" du site web ;
\task Donner aux élèves responsable du serveur web le fichier YAML et testez votre site marchand.
}
{[RES]_Gestion_des_produits}
{Moodle}

\end{document}
